\BeleskeID{02-s034}
% element: 02-s034:e01
% element: 02-s034:e02
% element: 02-s034:d01

Za jednu promenljivu prirodni raspored 0, 1 već zadovoljava uslov jediničnog Hemingovog rastojanja. Horizontalna i uspravna karta zato daju iste indekse: \(A=0\) odgovara polju 0, a \(A=1\) polju 1. Oznaka promenljive, vrednost iznad ili pored polja i indeks unutar polja imaju različite uloge.

Za dve promenljive niz 00, 01, 11, 10 daje indekse 0, 1, 3, 2. Preuređivanje je potrebno jer bi u prirodnom binarnom nizu prelaz 01 na 10 promenio dva bita. U Grejovom nizu susedni su i prvi i poslednji član: 00 i 10 razlikuju se u jednom bitu. Zamišljenim povezivanjem naspramnih ivica dobijamo prsten. U rasporedu \(2\times2\) iste kombinacije prepoznajemo iz oznaka reda i kolone.
